\section{Lelong numbers and divisorial estimates}
\label{sec:psh-divisorial-estimates}

Lelong numbers measure logarithmic concentration. Divisorial removal turns their lower bounds into inequalities of currents, while the local mass formula controls them under convergence and in a fixed background. For a prime divisor $E$ on a complex manifold, let $[E]$ be its integration current and $\nu(w,E)$ the generic Lelong number of a quasi-plurisubharmonic function $w$ along $E$.

We use the following precise local consequence of Siu's decomposition, as recalled in \cite{SGPT} and the standard current theory of \cite{DemAM11}. A closed positive $(1,1)$-current has, on a prime hypersurface $E$, coefficient-measure restriction equal to its generic Lelong number times $[E]$. In particular the remaining current after removal of any finite set of its divisorial components is positive. Distinct prime hypersurfaces do not contribute overlapping mass: their intersections have codimension at least two. These are statements about the local closed current.

\Needspace{8\baselineskip}
\begin{lemma}\label{lem:nonclosed-divisorial-removal}
Let $M$ be a complex manifold, let $\beta$ be a smooth real $(1,1)$-form on $M$, let $w\in\PSH(M,\beta)$, and let $E_1,\ldots,E_N$ be distinct prime divisors on $M$, with $N\geq1$. Then
\[
 \beta+\ddc w-\sum_{j=1}^N\nu(w,E_j)[E_j]\geq0.
\]
The form $\beta$ need not be closed.
\end{lemma}
\begin{proof}
Work on a coordinate ball and choose a smooth $\rho$ with $\ddc\rho\geq\beta$. The current
\[
 T=\ddc(w+\rho)=(\beta+\ddc w)+(\ddc\rho-\beta)
\]
is closed and positive. A smooth form, viewed through its coefficient measures, assigns zero mass to every proper analytic set. Thus $T$ and $\beta+\ddc w$ have the same restriction to $\bigcup_jE_j$. The stated local Siu fact identifies this restriction as $\sum_j\nu(w,E_j)[E_j]$. Subtracting it from the positive current $\beta+\ddc w$ leaves its coefficient-measure restriction to the complement, which is positive. This proves the inequality locally, hence globally. The argument does not identify $\beta+\ddc w$ itself as a closed current.
\end{proof}

\Needspace{12\baselineskip}
A potential whose curvature is a divisorial current plus a smooth form is determined, up to ordinary domination, by its divisorial coefficients. The comparison extends to countably many components of finite total mass.

\begin{lemma}\label{lem:divisorial-potential-domination}
Let $Y$ be a connected compact complex manifold. Let $D_j$ be distinct prime divisors, $c_j>0$ real coefficients such that $S=\sum_jc_j[D_j]$ is a positive current of finite mass, and $\beta$ a smooth real $(1,1)$-form; $[D_j]$ denotes the current of integration. Let $v\in\QPSH(Y)$ satisfy
\[
 \ddc v=S-\beta.
\]
For $a\in\QPSH(Y)$, the following conditions are equivalent:
\[
 a\preceq v,\qquad a\preceq_{\mathcal I}v,\qquad
 \nu(a,D_j)\geq c_j\quad\hbox{for every }j.
\]
\end{lemma}
\begin{proof}
The coefficient-measure restriction of $S$ to $D_j$ is $c_j[D_j]$. Distinct divisors have intersections of codimension at least two, which carry no mass of either divisor current. Countable summation preserves this assertion. Smooth forms give every divisor zero coefficient mass. Local Siu decomposition after a smooth correction therefore gives $\nu(v,D_j)=c_j$.

Ordinary domination implies ideal domination by integrability, and ideal domination implies the asserted divisorial Lelong inequalities by \cref{thm:local-i-inputs}. It remains to prove that these inequalities imply ordinary domination.

Choose a Hermitian form $\Lambda$ with $a\in\PSH(Y,\Lambda)$. By \cref{lem:nonclosed-divisorial-removal}, the divisorial inequalities allow removal of each finite partial sum of $S$ from $\Lambda+\ddc a$ while retaining positivity. Passing to the increasing finite sums gives
\[
 \Lambda+\ddc a-S\geq0.
\]
Consequently $\ddc(a-v)\geq\beta-\Lambda$. The $L^1$ difference has a quasi-plurisubharmonic representative and is bounded above on compact $Y$. This yields $a\leq v+C$ almost everywhere. \Cref{lem:qpsh-almost-everywhere-order} extends this inequality everywhere. This proves ordinary domination.
\end{proof}


\begin{lemma}\label{lem:lelong-lone-semicontinuity}
Let $u_j,u$ be quasi-plurisubharmonic on a complex manifold, with a common smooth lower curvature bound. If $u_j\to u$ in $L^1_{\mathrm{loc}}$, then at every point $x$,
\[
 \nu(u,x)\geq\limsup_j\nu(u_j,x).
\]
\end{lemma}
\begin{proof}
On a coordinate ball about $x$, add one smooth function to all the potentials to make them plurisubharmonic. Write the resulting potentials as $p_j,p$, let $n$ be the dimension of the chart, and let $\kappa$ be its Euclidean form. Distributional convergence gives weak convergence of the positive trace measures $\ddc p_j\wedge\kappa^{n-1}$ to $\ddc p\wedge\kappa^{n-1}$. For every sufficiently small radius $r$ whose boundary sphere has zero mass for the latter measure, the masses of the balls converge. The Lelong monotonicity formula bounds each $\nu(p_j,x)$ by its normalized trace mass on that ball. Taking the upper limit in $j$, and then letting such radii decrease to zero, gives $\limsup_j\nu(p_j,x)\leq\nu(p,x)$. Smooth additions do not change Lelong numbers.
\end{proof}

\begin{lemma}\label{lem:gauduchon-uniform-lelong}
Let $Z$ be connected compact complex and let $\beta$ be a smooth real $(1,1)$-form with $\PSH(Z,\beta)\ne\varnothing$. There is a finite constant $C_\beta$ such that
\[
 \nu(v,z)\leq C_\beta\qquad
 \text{for every }v\in\PSH(Z,\beta)\text{ and every }z\in Z.
\]
In particular all generic divisorial Lelong numbers in this fixed background are uniformly bounded. Additive normalization of the potentials is unnecessary.
\end{lemma}
\begin{proof}
Subtract the supremum from $v$. By \cref{thm:normalized-compactness}, these normalized potentials have a common $L^1$ bound. Cover $Z$ by finitely many nested coordinate balls. On each larger ball choose a fixed smooth quadratic function $r$ with $\ddc r\geq\beta$, so $p=v+r$ is plurisubharmonic there. The functions $p$ have uniformly bounded local $L^1$ norms.

Write $d=\dim Z$ and let $\kappa$ be the Euclidean form in the chart. Choose a nonnegative smooth cutoff $\chi$ supported in the larger ball and equal to one on an intermediate ball. Positivity and distributional integration by parts give
\[
 \int_{\text{intermediate ball}}\ddc p\wedge\kappa^{d-1}
 \leq\int\chi\,\ddc p\wedge\kappa^{d-1}
 =\int p\,\ddc\chi\wedge\kappa^{d-1}
 \leq C\|p\|_{L^1}.
\]
Every point of the smaller ball admits a surrounding Euclidean ball of a fixed radius inside the intermediate ball. The local Lelong monotonicity formula, recalled in \cite{SGPT}, bounds the Lelong number at its center by a fixed multiple of this trace mass. The finite cover gives one bound on $Z$. Smooth additions and additive normalization do not change Lelong numbers, proving the assertion for the original potentials.
\end{proof}


The Bergman estimates also control Lelong numbers on a fixed local modification. A \emph{local divisor over $X$} means a prime divisor $E$ on a smooth modification $\pi:Z\to U$ of a coordinate neighborhood. For a quasi-plurisubharmonic function $u$ on $X$, let $\nu(u\circ\pi,E)$ be its generic Lelong number along $E$, let $\operatorname{ord}_E$ denote divisorial vanishing order, and let $\operatorname{Jac}\pi$ be the local holomorphic Jacobian determinant. Write
\[
 \nu_E(u)=\nu(u\circ\pi,E),\qquad
 A_U(E)=1+\operatorname{ord}_E(\operatorname{Jac}\pi).
\]
The second number is the log discrepancy. These local modifications are not required to extend over the compact manifold.

\begin{lemma}\label{lem:bergman-divisorial-error}
For a quasi-plurisubharmonic function $u$, every fixed local divisor $E$, and every real $m>0$, one has
\[
 \max\!\left\{\nu_E(u)-\frac{A_U(E)}m,0\right\}
 \leq \frac1m\operatorname{ord}_E\mathcal I(mu)
 \leq \nu_E(u).
\]
Consequently, a local logarithmic presentation $a_m=m^{-1}\log\sum_k|f_{m,k}|^2+O(1)$, where the $f_{m,k}$ generate $\mathcal I(mu)$, satisfies $\nu_E(a_m)\to\nu_E(u)$ as $m\to\infty$.
\end{lemma}
\begin{proof}
The upper estimate follows from the local Bergman lower bound in \cref{prop:bergman-local-estimates}. The logarithm of local generators of $\mathcal I(mu)$, divided by $m$, is at least $u$ up to a locally bounded term. Pulling back and taking a generic Lelong number reverses this inequality.

For the lower estimate, take a local generator $f$ of the ideal and set $k=\operatorname{ord}_E(f)$. At a generic smooth point of $E$, choose a coordinate $z_1$ defining it. Divisorial removal, after a smooth correction, gives
\[
 u\circ\pi\leq\nu_E(u)\log|z_1|^2+C.
\]
\Needspace{9\baselineskip}
Shrink around a point where the residual factors of $f\circ\pi$ and of the Jacobian are nonvanishing. The change-of-variables formula and local integrability of $|f|^2\mathrm{e}^{-mu}$ then require integrability in $z_1$ of
\[
 |z_1|^{2(k+A_U(E)-1-m\nu_E(u))}.
\]
This implies $k+A_U(E)>m\nu_E(u)$. In particular the weaker non-strict lower bound in the statement holds. Take the minimum over a finite set of local generators and use nonnegativity of the order. Since $\nu_E(a_m)=m^{-1}\operatorname{ord}_E\mathcal I(mu)$, the last assertion follows by letting $m\to\infty$.
\end{proof}

For rational $m>0$, the presentation $a_m$ has analytic singularities in the convention of the source book. The real-parameter estimate does not require this restriction.


\begin{lemma}\label{lem:analytic-ideal-dominance}
Let $v$ be quasi-plurisubharmonic on a compact complex manifold, and let $a$ have analytic singularities. If $v\preceq_{\mathcal I}a$, then $v\preceq a$.
\end{lemma}
\begin{proof}
Resolve the coherent ideal of an analytic presentation of $a$ by a smooth modification $\pi:Y\to X$. In local coordinates $z_1,\ldots,z_n$ on $Y$, its expression has coefficients $b_i\geq0$, $1\leq i\leq r\leq n$, and a locally bounded remainder $B$:
\[
 a\circ\pi=\sum_{i=1}^r b_i\log|z_i|^2+B,
 \qquad b_i\geq0,\qquad B\in L^\infty_{\mathrm{loc}}.
\]
By \cref{thm:local-i-inputs}, the generic Lelong number of $v\circ\pi$ along $\{z_i=0\}$ is at least $b_i$. Add a smooth local function to $v\circ\pi$ to make it plurisubharmonic. Divisorial removal for this closed positive complex Hessian shows that
\[
 v\circ\pi-\sum_i b_i\log|z_i|^2
\]
has a quasi-plurisubharmonic representative. Here the difference is first interpreted off the divisor, or as an integrable distribution. Such a representative is locally bounded above. Since $B$ is locally bounded, $v\circ\pi\leq a\circ\pi+C$ on each relatively compact resolution chart away from the divisor. A finite cover above a smaller coordinate neighborhood makes $C$ uniform there. Push the inequality down off the analytic exceptional set; the canonical quasi-plurisubharmonic representatives extend it everywhere. Finitely many neighborhoods on the compact manifold give $v\leq a+C$ globally.

The only resolution input is the standard principalization of an analytic ideal. Divisorial removal is applied to a locally closed positive Hessian after smooth correction, not to a purported closed background on the original manifold.
\end{proof}


\Needspace{12\baselineskip}
\begin{lemma}\label{lem:analytic-minorant-product-ideals}
Let $u,v$ be quasi-plurisubharmonic on a compact complex manifold, and let $s$ have analytic singularities. If $u+v\leq s+O(1)$, then for analytic potentials $U_m,V_m$ of types $(1/m,\mathcal I(mu))$, $(1/m,\mathcal I(mv))$, there are numbers $\delta_m\downarrow0$ such that
\[
 U_m+V_m\leq(1-\delta_m)s+O(1)
 \quad\hbox{for all sufficiently large integers }m.
\]
The bounded error may depend on $m$.
\end{lemma}
\begin{proof}
Resolve the ideal of $s$ by a smooth modification $\pi:\widetilde X\to X$. Let $E_i$ be the components of its resolved divisor with positive coefficients $c_i$. In resolution coordinates $z_i$ defining these components, its singular part is $\sum_i c_i\log|z_i|^2$. Set
$A_i=1+\operatorname{ord}_{E_i}(\operatorname{Jac}\pi)$.
The hypothesis and \cref{lem:bergman-divisorial-error} give
\[
 \nu(U_m\circ\pi,E_i)+\nu(V_m\circ\pi,E_i)
 \geq\nu(u\circ\pi,E_i)+\nu(v\circ\pi,E_i)-\frac{2A_i}{m}
 \geq c_i-\frac{2A_i}{m}.
\]
Choose a positive integer $C\geq\max_i2A_i/c_i$ and put $\delta_m=C/m<1$ for large $m$. After a smooth local correction, divisorial removal from the positive Hessian of $(U_m+V_m)\circ\pi$ shows that
\[
 (U_m+V_m)\circ\pi-(1-\delta_m)\sum_i c_i\log|z_i|^2
\]
has a quasi-plurisubharmonic representative. It is locally bounded above. A finite resolution cover and proper descent therefore give
$U_m+V_m\leq(1-\delta_m)s+O(1)$, just as in \cref{lem:analytic-ideal-dominance}. The bounded constant may depend on $m$. If the divisor is empty, $s$ is bounded and upper bounds for $U_m,V_m$ prove the assertion directly.
\end{proof}
